\clearpage\begingroup
\ArithmeticHead{24}{L}{PHILIPP MAENNCHEN, GAUSS ALS ZAHLENRECHNER.}
\renewcommand{\thefootnote}{\arabic{footnote})}\setcounter{footnote}{0}
\begin{center}V.\\[3pt]{\large\textls[100]{Berechnung von reziproken Werten.}}\end{center}
\noindent\hspace*{\parindent}In dem zuletzt vorgeführten Beispiel hat die Zahl selbst — oder vielmehr ein 6stelliges Stück derselben — sich zur Faktorenzerlegung geeignet;
beim vorhergehenden Beispiel musste \textsc{Gauss} den reziproken Wert nehmen,
der allerdings in diesem Falle leicht zu finden war. Ob er es andernfalls
auch getan hätte? Nun, wir finden Beispiele, wo er auch unter erschwerenden
Umständen reziproke Werte berechnete. Auch hierin hatte er sich eine bedeutende Gewandtheit erworben, und wiederum war es der Einfluss der Empirie, der eine solche Übung notwendig machte. Wenn eine wirklich oder
vermeintlich wichtige Konstante ermittelt war, so war es möglich, dass ihr
Quadrat, ihre Quadratwurzel, ihr Logarithmus, ihr reziproker Wert oder dergleichen mit andern bereits ermittelten Konstanten einen erkennbaren
\clearpage\SourceLeaf{580}\ArithmeticHead{25}{R}{V. BERECHNUNG VON REZIPROKEN WERTEN.}
\noindent Zusammenhang hatte. Das konnte dann zu neuen Sätzen und Untersuchungen
führen.

Bei \textsc{Gauss} kam nun noch eine ganz besondere Veranlassung hinzu. Seine
Studien über das arithmetisch-geometrische Mittel führten ihn u. a. auch zu
dem Versuch, das zu bestimmende Mittel durch eine Reihenentwicklung zu
gewinnen. Es gelang ihm auch, $M(1+x,\;1-x)$ als eine Potenzreihe von $x$
darzustellen, allein die Koeffizienten von $x$ in dieser Entwicklung befolgten
kein leicht erkennbares Gesetz. Erst als er auf den Gedanken kam, den reziproken Wert, also $\frac{1}{M(1+x,\;1-x)}$ in eine Reihe zu entwickeln, gelangte er zu
dem schönen Ergebnis:
\[
\frac{1}{M(1+x,\;1-x)}=1+\frac14x^2+\frac9{64}x^4+\frac{25}{256}x^6+\frac{1225}{16384}x^8+\cdots.
\]
\noindent Die Koeffizienten der aufeinanderfolgenden Glieder sind nun die Quadrate
von $\frac12$, $\frac12\mathord{\cdot}\frac34$, $\frac12\mathord{\cdot}\frac34\mathord{\cdot}\frac56$, $\frac12\mathord{\cdot}\frac34\mathord{\cdot}\frac56\mathord{\cdot}\frac78$ usw.\footnote{Werke III, S. 366 u. ff.}.

Dieses Bildungsgesetz hat er wieder zunächst induktiv erschlossen; aber
er bleibt nicht dabei stehen, sondern er führt sofort den strengen Beweis für
seine Gültigkeit. Andererseits macht er aber auch Übungen in der zahlenmässigen Bestimmung von arithmetisch-geometrischen Mitteln nach dem neugefundenen Verfahren, und, um sie mit dem entsprechenden, nach dem seitherigen Verfahren ermittelten Wert vergleichen zu können, muss er von einem
der beiden Ergebnisse den reziproken Wert bestimmen.

Auf einem Blatte des Nachlasses\footnote{In Kapsel 50, Elliptische Funktionen, Fi, Nr. 5.} findet sich ein vollständig durchgeführtes Beispiel:
\[
\begin{array}{rr}
\multicolumn{2}{c}{M\sqrt{0{,}96}\text{\textsuperscript{3)}}}\\[3pt]
1{,}00000\,00000&0{,}00000\,00000\\
0{,}97979\,58971&9{,}99113\,56165\\
0{,}98984\,64001&9{,}99559\,04242\rlap{\text{\textsuperscript{4)}}}\\
0{,}98989\,79485&9{,}99556\,78082\\
0{,}98987\,21743&9{,}99557\,91162.
\end{array}
\]
\setcounter{footnote}{3}\footnotetext{$M\sqrt{0{,}96}$ bedeutet $M(1,\;\sqrt{0{,}96})$, und dies ist entstanden aus $M(1+\frac15,\;1-\frac15)$; die Variable $x$ hat also hier den Wert $\frac15$.}
\setcounter{footnote}{4}\footnotetext{In der Handschrift hatte \textsc{Gauss} die Logarithmen der zweiten und dritten Zahl mit einander vertauscht und das durch ein Zeichen angedeutet.}
\par\noindent X 2 Abh. 6.\hfill 4
\clearpage\SourceLeaf{581}\ArithmeticHead{26}{L}{PHILIPP MAENNCHEN, GAUSS ALS ZAHLENRECHNER.}
\noindent\hspace*{\parindent}Calculus posterior confirmatur adjumento seriei
\[
M\cos\varphi=\frac{1}{1+\frac14xx+\frac{1\mathord{.}9}{4\mathord{.}16}\mathord{\cdot}x^4+\text{etc.}}
\]
\begingroup\small\setlength{\arraycolsep}{3pt}
\[
\begin{array}{rl}
1{,}00000\,00000\mathord{.}00&\\
0{,}01000\,00000\mathord{.}00&\\
\mathord{.}\;\mathord{.}\;\mathord{.}\;22\,50000\mathord{.}00&\\
\mathord{.}\;\mathord{.}\;\mathord{.}\;\mathord{.}\;\mathord{.}\;62500\mathord{.}00&\\
\mathord{.}\;\mathord{.}\;1914\mathord{.}06&\\
\mathord{.}\;\mathord{.}\;\mathord{.}\;62\mathord{.}02&\\
\mathord{.}\;\mathord{.}\;\mathord{.}\;\mathord{.}\;2\mathord{.}08&\\
\mathord{.}7&\\\cline{1-1}
1{,}01023\,14478\mathord{.}23&\times0{,}99\\
1010\,23144\;78&\\\cline{1-1}
1{,}00012\,91333\;45&\times0{,}9999\\
10\,00129\;13&\\\cline{1-1}
1{,}00002\,91204\;32&\text{Rec.}\\[3pt]
0{,}99997\,08804\;16&\left[\frac{1}{1+\varepsilon}=1-\varepsilon+\varepsilon^2-+\cdots\right]\\[3pt]
9\,99970\;88&[\times\,0{,}99\ \text{wie oben}]\\\cline{1-1}
0{,}99987\,08833\;28&[\times\,0{,}9999\ \text{wie oben}]\\
999\,87088\;33&\\\cline{1-1}
0{,}98987\,21744\mathord{.}95.&
\end{array}
\]
\endgroup
\par\noindent\hspace*{\parindent}Dies ist nun der gesuchte reziproke Wert, der mit dem nach der ersten
Methode gefundenen auf 9 Stellen übereinstimmt. Man sieht, es wird zunächst durch geeignete Multiplikation, und zwar wieder unter geschickter Verwendung der welschen Praktik, eine Zahl hergestellt, die nur um einen sehr
kleinen Betrag $\varepsilon$ grösser ist als 1. Dann ist der reziproke Wert sehr nahe
gleich $1-\varepsilon+\varepsilon^2$, während die folgenden Glieder der unendlichen Reihe vernachlässigt werden können. Mit der sich ergebenden Zahl $1-\varepsilon+\varepsilon^2$ werden
nun die gleichen Multiplikationen vorgenommen wie mit der ursprünglichen
Zahl, und so erhält man alsdann den gesuchten reziproken Wert.

Eine andere Rechnung dieser Art, dazu mit grösserem Kraftaufwand,
aber ohne jede erläuternde Bemerkung, findet sich auf einem Zettel des
\clearpage\SourceLeaf{582}\ArithmeticHead{27}{R}{VI. GAUSS’ METHODEN, DIE FAKTOREN GROSSER ZAHLEN ZU FINDEN.}
\setcounter{footnote}{0}
\noindent Nachlasses\footnote{Arithmetik a 5, Kapsel 40.}. Die Bemerkungen zum vorigen Beispiel erleichtern das Verständnis
dieser Rechnung, die ein Torso geblieben ist.
\begingroup\footnotesize\setlength{\arraycolsep}{2pt}
\[
\begin{array}{rl}
603[{,}]5533905\,9327376220\,042218105&[=x]\\
1207\;1067811\,8654752440\,084436211&[=2x]\\
18\;1066017\,1779821286\,601266543&[=0{,}03x]\\
1189\;0001794\,6874931153\,483169668&[=1{,}97x=2x-0{,}03x]\\
41\;0000061\,8857756246\,671833436\;83&[:29]\\
1{,}\;0000001\,5094091615\,772483742\;3617&[:41]\\
1\,5000002264\,113742365\,8725&[\times15\mathord{.}10^{-8}]\\\cline{1-1}
1{,}\;0000000\,0094089351\,658741376\;4892.&
\end{array}
\]
\endgroup
\noindent Es würde zu sehr das Eingehen auf Einzelheiten erfordern, wenn wir bei den
interessanten Gedankengängen verweilen wollten, die dieser Rechnung zugrunde
liegen; es kann daher unter nochmaligem Hinweis auf das vorhergehende
Beispiel dem Leser überlassen werden\footnote{Auf $29\mathord{.}41=1189$ kommen wir an einer späteren Stelle noch einmal zu sprechen. Man beachte, dass \textsc{Gauss} hier durch 29 und 41 dividiert, ohne Reste aufzuschreiben.}. \textsc{Gauss} hat hier abgebrochen, wahrscheinlich, weil er sich bei der Multiplikation $15\mathord{.}10^{-8}$ verrechnet hat (statt
$1500000226\ldots$ steht nämlich in der Handschrift $150000225141137423658\ldots$).
Durch irgend eine Kontrollrechnung scheint er den Fehler bemerkt zu haben,
und so unterliess er es, den reziproken Wert zu ermitteln.
\clearpage\endgroup
